\section{Implementation}
We implement R2CC as a drop-in extension of 3K lines of C++ code to NCCL 2.23.4, allowing existing ML frameworks to leverage it without any modifications. Our implementation integrates R2CC strategies into the host transport layer and collective scheduling logic, and builds R2CC-AllReduce from standard NCCL collectives on sub-communicators.

\smallskip\noindent\textbf{Transport hooks and backup paths.}
R2CC intercepts NCCL's communicator initialization at the \texttt{ncclNet} plugin layer to establish backup connections alongside the primary connection. We extend the memory registration path to register each GPU buffer with all NICs of the server, and augment the connection setup logic to create additional queue pairs over backup NICs. The resulting handles are organized into a per-channel failover list ordered by PCIe distance to the source GPU, allowing the failure recovery logic to select the topologically closest healthy NIC.

\smallskip\noindent\textbf{Failure detection and live migration.}
R2CC augments the net plugin and proxy send receive progress loops to handle network failures without crashing the application process. The modified InfiniBand plugin monitors Completion Queue (CQ) and Queue Pair (QP) errors, reporting them as failed requests to the proxy. Upon error, the plugin issues a lightweight RDMA probe and exchanges messages via the bootstrap network to confirm the failure location. Once a failure is confirmed, the proxy rolls the channel back to the last chunk with a confirmed completion, purges outstanding work requests, transitions the channel to the backup connection, and reissues the remaining chunks.

\smallskip\noindent\textbf{R2CC-Balance and R2CC-AllReduce.}
R2CC-Balance is integrated into NCCL's enqueue and planning logic. Before NCCL selects an algorithm and instantiates rings or trees, R2CC inspects the health status records of NICs on each node. The planner then redistributes the affected channels across the remaining healthy NICs on the same node. R2CC-AllReduce creates two sub-communicators with \texttt{ncclCommSplit} after repair. Most of the data runs a global AllReduce across all nodes, while the rest runs a Partial AllReduce on the sub-communicator that excludes the degraded node. The degraded node's share of that part is then reduced onto a healthy helper rank and broadcast back in pipelined chunks. These collectives run on separate CUDA streams ordered by events, so the phases overlap without custom kernels. HotRepair, R2CC-Balance, and R2CC-AllReduce all support CUDA Graphs (Appendix~\ref{app:discussion}).
